% GENERATED FILE. Edit canonical lessons, then run npm run generate:books.
% canonical-source-hash: fnv1a64:643aaf8c2ec51424
% canonical-lessons: GU-C306-svayansevak, GU-C306-rahat, GU-C306-chavni, GU-C306-anaj, GU-C306-umatvu

\chapter{The Relief Camp}
\label{ch:gu-the-relief-camp}
\hlchaptermodality{306}

\begin{chapteropening}
\textbf{By the end of this chapter:} \emph{I can say who came to help and what they brought.}

Complete the last lesson of chapter 306: I can say who came to help and what they brought.
\end{chapteropening}

\section[\texorpdfstring{\gu{સ્વયંસેવક} (svayansevak)}{સ્વયંસેવક (svayansevak)}]{\gu{સ્વયંસેવક} (svayansevak) — a volunteer}

\label{lesson:GU-C306-svayansevak}



\begin{quote}
Before the new one: say the Gujarati for to be buried, then the Gujarati for rubble.
\end{quote}

\subsection*{You'll want to know: \gu{સ્વયંસેવક}}
\textbf{\gu{સ્વયંસેવક}} — \emph{svayansevak} — \textquotedblleft{}a volunteer\textquotedblright{}.

\begin{quote}
\gu{એક} \gu{સ્વયંસેવકે} \gu{બાળકોને} \gu{દૂધ} \gu{આપ્યું}.
\end{quote}

\begin{quote}
\emph{ek svayansevke bāḷkone dūdh āpyũ.}
\end{quote}

\begin{quote}
A volunteer gave the children milk.
\end{quote}

A masculine or feminine noun, by the person meant.

\begin{cousinweb}
\textbf{\gu{સ્વયંસેવક}} (\emph{svayansevak}) is built from Sanskrit \emph{svayam}, \textquotedblleft{}oneself, of one's own accord\textquotedblright{}, and \emph{sevaka}, \textquotedblleft{}one who serves\textquotedblright{} (Monier-Williams): someone who serves of their own will.
\end{cousinweb}

\begin{grammarlens}[title={What you've built}]
One more word for telling what happened.
\end{grammarlens}

\subsection*{Your turn}
\begin{itemize}
\raggedright
  \item \emph{Say it:} \emph{svayansevak}
  \item \emph{Say it:} \emph{svayansevak}, once more
  \item \emph{Say it:} \emph{ek svayansevke bāḷkone dūdh āpyũ.} — A volunteer gave the children milk.
  \item \emph{Recall:} say the Gujarati for to be buried, then the Gujarati for rubble, then say \emph{svayansevak} again
\end{itemize}

\begin{tcolorbox}[breakable,colback=teal!4,colframe=teal!35!black,title={Before you move on}]
What does \gu{સ્વયંસેવક} mean? (\textquotedblleft{}A volunteer\textquotedblright{}.) Say it once more.
\end{tcolorbox}
\section[\texorpdfstring{\gu{રાહત} (rāhat)}{રાહત (rāhat)}]{\gu{રાહત} (rāhat) — relief}

\label{lesson:GU-C306-rahat}



\begin{quote}
Before the new one: say the Gujarati for rubble, then the Gujarati for a volunteer.
\end{quote}

\subsection*{You'll want to know: \gu{રાહત}}
\textbf{\gu{રાહત}} — \emph{rāhat} — \textquotedblleft{}relief\textquotedblright{}.

\begin{quote}
\gu{વરસાદ} \gu{અટક્યો} \gu{ત્યારે} \gu{લોકોને} \gu{રાહત} \gu{થઈ}.
\end{quote}

\begin{quote}
\emph{varsād aṭakyo tyāre lokone rāhat thaī.}
\end{quote}

\begin{quote}
When the rain stopped, people felt relieved.
\end{quote}

A feminine noun.

\begin{cousinweb}
\textbf{\gu{રાહત}} (\emph{rāhat}) is usually traced, through Persian, to Arabic \emph{rāḥa}, \textquotedblleft{}rest, ease\textquotedblright{}. Gujarati uses it for the ease that comes after trouble, \textbf{\gu{રાહત} \gu{થઈ}}, and for the help sent to people in need.
\end{cousinweb}

\begin{grammarlens}[title={What you've built}]
One more word for telling what happened.
\end{grammarlens}

\subsection*{Your turn}
\begin{itemize}
\raggedright
  \item \emph{Say it:} \emph{rāhat}
  \item \emph{Say it:} \emph{rāhat}, once more
  \item \emph{Say it:} \emph{varsād aṭakyo tyāre lokone rāhat thaī.} — When the rain stopped, people felt relieved.
  \item \emph{Recall:} say the Gujarati for rubble, then the Gujarati for a volunteer, then say \emph{rāhat} again
\end{itemize}

\begin{tcolorbox}[breakable,colback=teal!4,colframe=teal!35!black,title={Before you move on}]
What does \gu{રાહત} mean? (\textquotedblleft{}Relief\textquotedblright{}.) Say it once more.
\end{tcolorbox}
\section[\texorpdfstring{\gu{છાવણી} (chāvṇī)}{છાવણી (chāvṇī)}]{\gu{છાવણી} (chāvṇī) — a camp}

\label{lesson:GU-C306-chavni}



\begin{quote}
Before the new one: say the Gujarati for a volunteer, then the Gujarati for relief.
\end{quote}

\subsection*{You'll want to know: \gu{છાવણી}}
\textbf{\gu{છાવણી}} — \emph{chāvṇī} — \textquotedblleft{}a camp\textquotedblright{}.

\begin{quote}
\gu{ગામના} \gu{લોકો} \gu{હવે} \gu{એક} \gu{છાવણીમાં} \gu{રહે} \gu{છે}.
\end{quote}

\begin{quote}
\emph{gāmnā loko have ek chāvṇīmā̃ rahe che.}
\end{quote}

\begin{quote}
The villagers are living in a camp now.
\end{quote}

A feminine noun.

\begin{cousinweb}
\textbf{\gu{છાવણી}} (\emph{chāvṇī}) goes back to Sanskrit \emph{chādana}, \textquotedblleft{}covering\textquotedblright{}, and CDIAL lists the Gujarati noun, \textquotedblleft{}huts, a cantonment\textquotedblright{} (CDIAL 5017): a place of shelters put up for a time.
\end{cousinweb}

\begin{grammarlens}[title={What you've built}]
One more word for telling what happened.
\end{grammarlens}

\subsection*{Your turn}
\begin{itemize}
\raggedright
  \item \emph{Say it:} \emph{chāvṇī}
  \item \emph{Say it:} \emph{chāvṇī}, once more
  \item \emph{Say it:} \emph{gāmnā loko have ek chāvṇīmā̃ rahe che.} — The villagers are living in a camp now.
  \item \emph{Recall:} say the Gujarati for a volunteer, then the Gujarati for relief, then say \emph{chāvṇī} again
\end{itemize}

\begin{tcolorbox}[breakable,colback=teal!4,colframe=teal!35!black,title={Before you move on}]
What does \gu{છાવણી} mean? (\textquotedblleft{}A camp\textquotedblright{}.) Say it once more.
\end{tcolorbox}
\section[\texorpdfstring{\gu{અનાજ} (anāj)}{અનાજ (anāj)}]{\gu{અનાજ} (anāj) — grain (as food)}

\label{lesson:GU-C306-anaj}



\begin{quote}
Before the new one: say the Gujarati for relief, then the Gujarati for a camp.
\end{quote}

\subsection*{You'll want to know: \gu{અનાજ}}
\textbf{\gu{અનાજ}} — \emph{anāj} — \textquotedblleft{}grain (as food)\textquotedblright{}.

\begin{quote}
\gu{છાવણીમાં} \gu{રોજ} \gu{અનાજ} \gu{અને} \gu{દૂધ} \gu{આવે} \gu{છે}.
\end{quote}

\begin{quote}
\emph{chāvṇīmā̃ roj anāj ane dūdh āve che.}
\end{quote}

\begin{quote}
Grain and milk come to the camp every day.
\end{quote}

A neuter noun.

\begin{cousinweb}
goes back to Sanskrit \emph{annādya}, \textquotedblleft{}food\textquotedblright{}, and CDIAL lists the Gujarati noun, \textquotedblleft{}grain\textquotedblright{} (CDIAL 398). Wheat, rice and millet are all \textbf{\gu{અનાજ}}.
\end{cousinweb}

\begin{grammarlens}[title={What you've built}]
One more word for telling what happened.
\end{grammarlens}

\subsection*{Your turn}
\begin{itemize}
\raggedright
  \item \emph{Say it:} \emph{anāj}
  \item \emph{Say it:} \emph{anāj}, once more
  \item \emph{Say it:} \emph{chāvṇīmā̃ roj anāj ane dūdh āve che.} — Grain and milk come to the camp every day.
  \item \emph{Recall:} say the Gujarati for relief, then the Gujarati for a camp, then say \emph{anāj} again
\end{itemize}

\begin{tcolorbox}[breakable,colback=teal!4,colframe=teal!35!black,title={Before you move on}]
What does \gu{અનાજ} mean? (\textquotedblleft{}Grain (as food)\textquotedblright{}.) Say it once more.
\end{tcolorbox}
\section[\texorpdfstring{\gu{ઊમટવું} (ūmaṭvũ)}{ઊમટવું (ūmaṭvũ)}]{\gu{ઊમટવું} (ūmaṭvũ) — to pour in; to throng (of people)}

\label{lesson:GU-C306-umatvu}



\begin{quote}
Before the new one: say the Gujarati for a camp, then the Gujarati for grain.
\end{quote}

\subsection*{You'll want to know: \gu{ઊમટવું}}
\textbf{\gu{ઊમટવું}} — \emph{ūmaṭvũ} — \textquotedblleft{}to pour in; to throng (of people)\textquotedblright{}.

\begin{quote}
\gu{આખા} \gu{શહેરમાંથી} \gu{લોકો} \gu{મદદ} \gu{લઈને} \gu{ઊમટ્યા}.
\end{quote}

\begin{quote}
\emph{ākhā śahermā̃thī loko madad laīne ūmaṭyā.}
\end{quote}

\begin{quote}
People poured in from all over the city, bringing help.
\end{quote}

Past (it …): \gu{તે} \gu{ટોળું} \gu{ઊમટ્યું}, \emph{te ṭoḷũ ūmaṭyũ}.

\begin{cousinweb}
\textbf{\gu{ઊમટવું}} (\emph{ūmaṭvũ}) probably goes back to Sanskrit \emph{únmṛṣṭa}, \textquotedblleft{}rubbed off\textquotedblright{}, and so \textquotedblleft{}made clear\textquotedblright{}: CDIAL lists Gujarati \emph{umaṭvũ}, \textquotedblleft{}to appear, come to light\textquotedblright{}, there (CDIAL 2133). The sense of crowds that \textbf{\gu{ઊમટે} \gu{છે}}, pour out all at once, is the modern language's.
\end{cousinweb}

\begin{grammarlens}[title={What you've built}]
One more word for telling what happened.
\end{grammarlens}

\subsection*{Your turn}
\begin{itemize}
\raggedright
  \item \emph{Say it:} \emph{ūmaṭvũ}
  \item \emph{Say it:} \emph{ūmaṭvũ}, once more
  \item \emph{Say it:} \emph{ākhā śahermā̃thī loko madad laīne ūmaṭyā.} — People poured in from all over the city, bringing help.
  \item \emph{Recall:} say the Gujarati for a camp, then the Gujarati for grain, then say \emph{ūmaṭvũ} again
\end{itemize}

\begin{tcolorbox}[breakable,colback=teal!4,colframe=teal!35!black,title={Before you move on}]
What does \gu{ઊમટવું} mean? (\textquotedblleft{}To pour in; to throng (of people)\textquotedblright{}.) Say it once more.
\end{tcolorbox}
